%% parts/p12/rewriting.tex --- Part 12.4
\chapter{Term Rewriting}\label{sec:p12-rewriting}

\begin{hbdefinition}{def:p12-trs}{First-order term-rewriting system}
Fix a ranked signature \(\Sigma\) and variables \(V\). A rewrite rule is a
pair \(l\to r\) of finite \(\Sigma\)-terms such that \(l\notin V\) and
\(\operatorname{Var}(r)\subseteq\operatorname{Var}(l)\). A term-rewriting
system \(R\) is a set of such rules, and
\[
 C[l\sigma]\to_R C[r\sigma]
\]
for every one-hole context \(C[-]\) and substitution \(\sigma\). The exact
name is needed to distinguish this contextual first-order relation from graph,
higher-order, modulo-AC, or strategy-restricted rewriting. It stands as a
binary relation on terms, differs through the displayed closure conditions,
and is the relation used by every result in this chapter and on
\hyperref[sheet:term-rewriting]{\texttt{sheet:term-rewriting}}.
\index{term-rewriting system}\index{rewrite relation}
\end{hbdefinition}

Write \(\to_R^*\) for reflexive transitive closure. A term is a normal form
when it has no \(R\)-successor. The relation is \emph{terminating} when it has
no infinite forward chain, \emph{confluent} when
\(s\to_R^*t_1\) and \(s\to_R^*t_2\) imply that \(t_1,t_2\) have a common
descendant, and \emph{locally confluent} when the same holds for one-step
forks. These properties have different quantifier patterns and are not
interchangeable.

\begin{hbtheorem}{thm:p12-newman}{cited=Newman1942}{Newman's lemma}
Every terminating and locally confluent rewrite relation is confluent.
\end{hbtheorem}

\noindent\emph{Proof.}
Use well-founded induction on the source \(s\) under the strict transitive
closure of rewriting. For reductions
\(s\to^*t_1\) and \(s\to^*t_2\), the empty-side cases are immediate.
Otherwise take first steps \(s\to u_1\) and \(s\to u_2\). Local confluence
joins \(u_1,u_2\) at some \(v\). The induction hypothesis, applied below
\(s\), first joins each tail \(u_i\to^*t_i\) with \(u_i\to^*v\), and then
joins the two resulting descendants. Termination is used exactly to justify
the induction; local confluence alone is insufficient.

\begin{hbdefinition}{def:p12-critical-pair}{Critical overlap and critical pair}
Rename two rules \(l_1\to r_1\), \(l_2\to r_2\) apart. If a nonvariable
subterm of \(l_2\) unifies with \(l_1\) with most general unifier \(\sigma\),
the two corresponding one-step reducts of \(l_2\sigma\) form a critical pair.
The nonvariable and rename-apart conditions are necessary to exclude ordinary
substitution compatibility and accidental variable capture. A critical pair
stands as a finite overlap witness, differs from an arbitrary peak by being
most general, and is used in the critical-pair criterion
\cite{Huet1973,Huet1980}.
\index{critical pair}\index{overlap!term rewriting}
\end{hbdefinition}

\begin{hbtheorem}{thm:p12-critical-pair}{cited=Huet1980}{Critical-pair criterion}
For a left-linear first-order term-rewriting system, local confluence is
equivalent to joinability of every critical pair. Without left-linearity,
variable-overlap phenomena require the corresponding stronger formulation.
\end{hbtheorem}

\begin{hbexample}{def:p12-rewrite-example}{A terminating nonconfluent system}
Let \(R=\{a\to b,a\to c\}\) with distinct constants. The system terminates,
but the one-step peak \(b\leftarrow a\to c\) is not joinable, so it is not
locally confluent and hence not confluent. Adding \(b\to c\) joins the peak
and preserves termination under the order \(a\succ b\succ c\).
\end{hbexample}

\section{Reduction orderings}

\begin{hbdefinition}{def:p12-reduction-order}{Reduction ordering}
A reduction ordering \(\succ\) is a well-founded strict order on terms that is
closed under substitutions and contexts. Orienting every rule \(l\to r\) with
\(l\succ r\) proves termination. The term is needed because an arbitrary
well-founded order need not survive substitution or contextual insertion. It
stands as a metatheoretic order on terms, differs through those stability
conditions, and is used by completion and termination proofs
\cite{Dershowitz1982,BaaderNipkow1998}.
\index{reduction ordering}
\end{hbdefinition}

Lexicographic path order (LPO) is determined by a precedence on function
symbols and compares root symbols together with recursively smaller argument
tuples; it is particularly suited to recursive definitions. Knuth--Bendix
order (KBO) assigns positive symbol weights and a precedence, compares total
weight first, and applies variable-occurrence conditions before precedence and
lexicographic tie-breaking. Neither acronym identifies one order until the
precedence, weights, status, and variable conditions are supplied.

\section{Knuth--Bendix completion}

\begin{hbmethod}{def:p12-completion}{Completion state and invariant}
A completion state is \((E,R)\), where \(E\) is a finite set of unoriented
equations and \(R\) a set of rules oriented by one fixed reduction ordering.
The transformations Orient, Simplify, Compose, Collapse, and Deduce are exactly
those on \hyperref[sheet:term-rewriting]{\texttt{sheet:term-rewriting}}. The
invariant is
\[
 \equiv_{E_0}=\equiv_{E\cup R},
\]
where both sides denote the generated equational congruence. The method name
is necessary to distinguish this fair transformation process from arbitrary
rule orientation. It stands as a semidecision method, differs because it may
fail or diverge, and is used in the conditional correctness theorem
\cite{KnuthBendix1970}.
\index{Knuth--Bendix completion}
\end{hbmethod}

Each transformation preserves the generated congruence: orientation replaces
an equation by one direction while equational closure keeps symmetry; simplify
and compose replace a side by an equal reduct; collapse returns an invalidated
left side to the equation set; and deduce adds an equality already generated by
two rule steps.

\begin{hbtheorem}{thm:p12-completion-correctness}{cited=Huet1980}{Conditional correctness of fair completion}
If a fair completion run terminates successfully with \(E=\varnothing\), every
critical pair of \(R\) joinable, and every rule decreasing under the fixed
reduction ordering, then \(R\) is terminating and confluent and
\[
 s\equiv_{E_0}t
 \quad\Longleftrightarrow\quad
 \operatorname{nf}_R(s)=\operatorname{nf}_R(t).
\]
The theorem is conditional: a failed orientation or an infinite fair run is
not a negative answer to the existence of some convergent presentation.
\end{hbtheorem}

Rewriting modulo associativity/commutativity, conditional rewriting, and
higher-order rewriting change unification, overlaps, and ordering requirements.
The first-order critical-pair and completion results above do not transfer
without their respective extensions.
